\scanpage{129}
\mspage{64}
$\bar L$ au-dessus de $\tilde U_x$, bien sûr, et
\[
 (34)\qquad
 1\mathrel{\text{---}}T^{I(x)}\longrightarrow
 \mathcal E_{\tilde U_x}^{\bar L}\longrightarrow
 \mathcal E_{k(x)}^{\overline{k(x)}}\longrightarrow 1
\]
où \deleted{$I(x)$}
\[
 I(x)=\{i\in I\mid x\in D_i\},\qquad T=T_\infty(\bar L^*)
\]
$\overline{k(x)}$ étant la clôture alg. de $k(x)$ induite par un plongement choisi de $\tilde K_x$ (corps des fonctions de $\tilde O_x$) dans $\bar L$, qui fait donc fonction de clôture sép. de $K_x$.

Ceci dit, pour $V$ anneau de valuation de $L/K$, si son centre $x$ \inserted{sur $X$} \unsure{est} $\in\operatorname{Supp}D$, on aura, pour $S=\operatorname{Spec}V$, d'où $\tilde S$ \inserted{(hensélisé)} du pt \unsure{générique|géométrique} $\tilde\eta$ une factorisation de $\tilde\eta\longrightarrow U$ (composé $\tilde\eta\longrightarrow\eta=\operatorname{Spec}(L)\longrightarrow U$) en $\tilde\eta\longrightarrow\tilde U_x$, car $\tilde\eta\mathrel{\text{---}}U$ est induite par $\tilde S\longrightarrow X$ qui se factorise en $\tilde S\longrightarrow\tilde T\longrightarrow X$ ($\tilde T=\operatorname{Spec}\tilde O_x$). Donc \deleted{l'image $N_V$ de} \inserted{l'hom. induit} $\mathcal E_{\tilde\eta}^{L}\longrightarrow\mathcal E_U^{L}$ est \deleted{contenue dans celle de} \inserted{un composé}
\[
 (35)\qquad
 \diagram{\begin{tikzcd}[column sep=large,row sep=small]
 \mathcal E_{\tilde\eta}^{\bar L}\arrow[r]\arrow[d,equal] &
 \mathcal E_{\tilde U_x}^{\bar L}\arrow[r]\arrow[d,equal,"{\unsure{\mathrm{dfn}|\mathrm{déf.}}}"'] &
 \mathcal E_U^{\bar L}\\
 N_V & \deleted{\unsure{N_x}}\ N_x &
 \end{tikzcd}}
\]
(l'image de $\mathcal E_{\tilde U_x}^{\bar L}\longrightarrow\mathcal E_U^{\bar L}$ n'étant d'ailleurs autre que le groupe de décomposition de $x$ dans $\mathcal E_U^{\bar L}$ relatif à $\bar L$\ldots). L'étude de l'hom composé
\margin{128}
